\section{Future Work}\label{sec:future}
The main open question is whether the corpus effect in RQ1 and RQ4 comes from the text domain or from annotation conventions. We ran preliminary versions of steps (1)--(3) after seeing the main results; however, they were designed post hoc and are not yet validated, therefore we do not report them here.
(1)~\textit{Error analysis:} every cross-corpus error is classified as a missed entity, a boundary error, a type error or a spurious prediction, and strict matching is compared with relaxed (overlap) matching and per-type F1, to see whether the gap reflects recognition failures or different definitions of an entity.
(2)~\textit{Guideline-specific heads:} one encoder is trained jointly on a language's two corpora with a separate adapter and output head for each guideline~\cite{stymne2018,crichton2017,houlsby2019,pfeiffer2020}. Scoring the same test text through each head then estimates how much of the cross-corpus gap the guideline alone accounts for.
(3)~\textit{Relabelling a related source:} titles, role nouns and generic institution nouns in the Sepedi and Sesotho training labels are rewritten towards news conventions, using a word list checked by at least two Setswana speakers per word, and the result is compared with a matched random-relabelling control. If a relabelled Sotho-Tswana source overtakes Swahili, the RQ4 reversal is a labelling effect.
(4)~\textit{A causal test:} two Setswana annotators relabel a blind sample of 200 test sentences from both corpora under one shared guideline (about 10--12 person-hours), and the cross-corpus gap is recomputed from the saved predictions.
(5)~\textit{Exposure and uncertainty:} an XLM-R-base contrast, which shares AfroXLMR's tokenizer but not its exposure, is the proper test for RQ2~\cite{conneau2020,alabi2022}. Moreover, more seeds, confidence intervals and crossed source and target random effects would tighten every estimate.
